\section{Conclusion}\label{sec:conclusion}

We set out to make a parallel timed-automata reachability engine faster, assuming
the obstacle would be work distribution.  It was not.  The engine expands at
$99.8\%$ parallel efficiency with milliseconds of cumulative load imbalance, and
the OpenMP chunk size---the knob one reaches for first---moves end-to-end runtime
by under $3\%$.  What limited it were two costs that parallelism cannot touch
and, worse, actively conceals.

The first was an Amdahl term hiding behind a latched trigger: a reclamation pass
guarded by a monotone counter, which turned an occasional serial sweep into a
per-level one whose cost grows with the state space.  The second, and the one we
think generalises furthest, is the \emph{compression tax}.  Minimal constraint
systems are the right way to store zones and cut our passed-set memory by roughly
$12\times$; but the zone algebra consumes dense closed matrices, so every
subsumption query rebuilt one at $O(|\clocks|^3)$.  With $89.8\%$ of insertions
being duplicates rejected by the first test, against antichains of mean length
$1.00$, we spent $22\times$ longer decompressing stored zones than comparing
against them.

The fix is to stop treating the compressed form as a blob to be restored and
start treating it as what it is---a generating set.  Inclusion against a minimal
constraint system is decidable in $O(|G|)$ directly (\cref{lem:minclosure}), and
because plain inclusion implies the coarser $\preceq_{aLU}$ subsumption the
engine actually prunes with (\cref{lem:filter}), the cheap exact test composes as
a pre-filter that provably preserves the explored state set
(\cref{thm:samestates}).  Forty lines, no new abstraction, no soundness debt.

If there is one transferable lesson, it is that \emph{a compressed representation
must be chosen with its query pattern in view, not only its size}.  The memory
win is easy to measure and gets measured; the per-query toll is paid on a hot
path, shows up in no benchmark, and---on a parallel tool---is masked by the fact
that the tool still scales beautifully.  It simply scales from a worse constant.
We suspect this pattern is not unique to us, nor to timed automata.

\paragraph{What we could not fix.}
We report two negative results, because they cost us more time than the positive
ones and may save a reader some.

\emph{Host-side bucketing in the GPU backend does not want to be parallel.}  A
shared lock-free table was measured against the serial loop and lost; we then
tried thread-private maps with a parallel counting-sort scatter, which avoids
that failure mode entirely, and gained ${\approx}1\%$ of wall-clock.  The loop is
bound by streaming the expansion output out of unified memory while the GPU is
using the same bandwidth, not by hashing.

\emph{The GPU backend cannot be routed by a simple occupancy heuristic.}
\Cref{sec:eval} shows viking15 and fischer11 differ $2.9\times$ in per-item GPU
throughput, tracking dispatch size---but arbiter17 has the largest dispatches of
all and still loses to the CPU, because its host and merge path costs
\SI{2.97}{\second} on top of a fast expansion.  Any threshold separating viking
from fischer misroutes arbiter.

\paragraph{Future work.}
The GPU backend stores zones densely and sizes its unified-memory pools to the
widest BFS level, costing $5$--$41\times$ the CPU backend's memory
(\cref{tab:mem}); porting \cref{sec:mingraph}'s representation to it, now that
\cref{sec:filter} makes such a representation cheap to query, is the obvious next
step and would attack both the memory and the host path at once.  Backend
selection deserves a genuine cost model spanning expansion \emph{and} merge
rather than a single statistic.  Finally, the order-sensitivity of convex-union
merging makes stored state counts vary run to run on some models
(\cref{sec:correctness}); a canonical merge order would restore replay
determinism, which matters more for debugging and artefact evaluation than for
verdicts.
